\section{Mixed endpoints and all-weight cyclotomic product quotients}
Put $A_{a,b}=\Imag\Li_{a,b}(\iu,\iu)$ and
$K_{a,b}=\Imag\Li_{a,b}(\iu,-1)$; retain $g_{a,b}$ and $S_p$ as above.
The all-exponent endpoint identity relates the mixed family to the same-color
family before any period-independence question is asked.

\subsection{An all-exponent mixed endpoint identity}

\begin{theorem}[Gaussian mixed endpoint]\label{cycloquot:thm:mixed-endpoint}
For every integer \(p\geq1\),
\begin{equation}\label{cycloquot:eq:mixed-endpoint}
 K_{p,1}+(-1)^{p-1}A_{1,p}
 =\beta(p+1)-2\log2\,\beta(p)
       +\sum_{j=2}^{p}(-1)^j\eta(j)\beta(p+1-j).
\end{equation}
Moreover,
\begin{equation}\label{cycloquot:eq:harmonic-split}
 S_p=g_{p,1}+K_{p,1}.
\end{equation}
\end{theorem}

\begin{proof}
All the boundary values used here converge. First-index-one values
with first argument \(\iu\) can be obtained by Dirichlet's test or
radial Abel limits. The shuffle and stuffle identities may therefore
be justified by convergent iterated integrals and Abel limiting.

Fix \(w=p+1\), and define homogeneous generating polynomials of
degree \(w-2\):
\[
 \mathcal A(X,Y)=\sum_{a+b=w}A_{a,b}X^{a-1}Y^{b-1},
 \quad
 \mathcal K(X,Y)=\sum_{a+b=w}K_{a,b}X^{a-1}Y^{b-1}.
\]
Let \(\mathcal V\) denote the same polynomial for
\(\Imag\Li_{a,b}(-1,\iu)\), and put
\[
 P(X,Y)=\sum_{a+b=w}
  \Imag[\Li_a(-1)\Li_b(\iu)]X^{a-1}Y^{b-1},
 \qquad
 H(X,Y)=\sum_{a+b=w}X^{a-1}Y^{b-1}.
\]
Stuffle and shuffle respectively give
\[
 \mathcal V(X,Y)+\mathcal K(Y,X)=P(X,Y)+\beta(w)H(X,Y),
\]
\[
 \mathcal A(X+Y,X)-\mathcal V(X+Y,Y)=P(X,Y).
\]
The collision term in the first equation is
\(-\Imag\Li_w(-\iu)=\beta(w)\).
Eliminating \(\mathcal V\) yields
\[
 \mathcal K(U,V)=P(V-U,U)+P(V,U)+\beta(w)H(V,U)
                  -\mathcal A(V,V-U).
\]
Set \(V=0\) and compare coefficients of \(U^{w-2}\).
The evaluation \(\Li_j(-1)=-\eta(j)\) gives
\eqref{cycloquot:eq:mixed-endpoint}, including the two occurrences of
\(-\log2\,\beta(p)\).

For \eqref{cycloquot:eq:harmonic-split}, taking imaginary parts of the outer
index in \(\Li_{p,1}(\iu,1)+\Li_{p,1}(\iu,-1)\) retains indices
\(2n+1\), with sign \((-1)^n\). Their inner sum is
\[
 \sum_{m=1}^{2n}\frac{1+(-1)^m}{m}=H_n.
\]
Absolute convergence proves the result directly for \(p>1\).
An Abel limit gives the endpoint \(p=1\).
\end{proof}

At \(p=4\) this specializes to
\begin{equation}\label{cycloquot:eq:K41}
 K_{4,1}=A_{1,4}+\frac{191\pi^5}{23040}
                  -\frac34G\zeta(3)-2\beta(4)\log2.
\end{equation}

\subsection{A precisely defined formal presentation}

Fix a level \(L=3\) or \(4\) and a weight \(w\geq2\). Let
\(\omega=e^{2\pi\iu/L}\).
Introduce rational formal symbols
\[
 D_a(r,s)\quad(1\leq a<w,\quad r,s\in\Z/L\Z)
\]
for imaginary parts of \(\Li_{a,w-a}(\omega^r,\omega^s)\).
Impose conjugation \(D_a(-r,-s)=-D_a(r,s)\).
Symbols fixed by this involution vanish. Put all single values and
imaginary parts of products of two single values in the background,
and also set \(D_a(1,0)=0\). The latter is precisely quotienting by
the chosen \(g\)-family.

The only further rows imposed are, for \(p+q=w\),
\begin{align}
 D_p(r,s)+D_q(s,r)&=0,\label{cycloquot:eq:formal-stuffle}\\
 \sum_{j=0}^{p-1}\binom{q+j-1}{j}D_{q+j}(s,r-s)
 +\sum_{j=0}^{q-1}\binom{p+j-1}{j}D_{p+j}(r,s-r)&=0.
 \label{cycloquot:eq:formal-shuffle}
\end{align}
They are the stuffle and shuffle rows of two single values, with their
one-leading-one regularizations. A regularized product involving the
real parameter \(T\) belongs to the background. At \(w=2\), the only
two-divergent-factor correction is real and disappears on taking
imaginary parts. Call the resulting quotient \(V_{L,w}\).

This is \emph{not} the full cyclotomic double-shuffle algebra. Products
involving a depth-two factor, relations lifted through higher depths,
additional functional equations, and motivic period assertions have
not been imposed. In particular, Zhao's standard relations
\cite{cycloquot:Zhao} are substantially broader than this presentation.
The theorem below computes \(V_{L,w}\), not a dimension of actual
numbers.

For a scalar assignment to the symbols, introduce
\[
 \mathcal F^{r,s}(X,Y)=
 \sum_{a=1}^{w-1}D_a(r,s)X^{a-1}Y^{w-a-1}.
\]
The conditions that the assignment annihilate all rows are exactly
\begin{align}
 \mathcal F^{r,s}(X,Y)+\mathcal F^{s,r}(Y,X)&=0,
 \label{cycloquot:eq:poly-stuffle}\\
 \mathcal F^{s,r-s}(X+Y,X)+
       \mathcal F^{r,s-r}(X+Y,Y)&=0.
 \label{cycloquot:eq:poly-shuffle}
\end{align}
These follow by comparing binomial coefficients after substitution.
They describe the dual of \(V_{L,w}\), so their solution dimension is
the required quotient dimension.

\subsection{All-weight dimension theorems}

\begin{theorem}[Level four]\label{cycloquot:thm:level-four}
The formal quotient satisfies
\[
 \dim_\Q V_{4,w}=
 \begin{cases}
 0,&w\ \text{even},\\
 (w-1)/2,&w\ \text{odd}.
 \end{cases}
 \qquad
 \sum_{w\geq2}\dim V_{4,w}\,t^w=\frac{t^3}{(1-t^2)^2}.
\]
At odd weight, its dual is parameterized by arbitrary antisymmetric
homogeneous polynomials of degree \(w-2\).
\end{theorem}

\begin{proof}
Up to conjugation, the six color pairs are
\((0,1),(1,0),(1,1),(1,2),(1,3),(2,1)\).
Write their polynomials as \(U,G,A,K,M,V\), respectively, and let
\(n=w-2\). The background condition gives \(G=0\);
stuffle gives
\[
 U=0,\quad A(Y,X)=-A(X,Y),\quad
 V(X,Y)=-K(Y,X),\quad M(Y,X)=M(X,Y).
\]
The shuffle with colors \((1,0)\) forces \(M=0\).
The shuffle with colors \((2,1)\) gives
\[
 A(X+Y,X)+K(Y,X+Y)=0,
\]
hence \(K(X,Y)=-A(Y,Y-X)\).
The shuffle with colors \((1,3)\) is
\(K(X,Y)=K(X,X-Y)\).
After substitution and antisymmetry this condition is
\(A=(-1)^{n+1}A\).

For even \(n\) all polynomials therefore vanish.
For odd \(n\), any antisymmetric \(A\), with \(K,V\) as above and
\(U=G=M=0\), satisfies the rows. To check exhaustiveness, the
unordered color classes up to conjugation are
\[
 (0,0),(0,1),(0,2),(1,1),(1,2),(1,3),(2,2).
\]
The real classes vanish, and the remaining classes are precisely the
conditions just used, together with their transposes. Thus there are
no unexamined rows.
A basis for \(A\) at odd \(w\) is
\[
 X^{a-1}Y^{w-a-1}-X^{w-a-1}Y^{a-1},
 \qquad 1\leq a\leq(w-1)/2.
\]
Counting these basis vectors and summing over odd weights proves
the formulas.
\end{proof}

\begin{theorem}[Level three]\label{cycloquot:thm:level-three}
The formal quotient satisfies
\[
 \dim_\Q V_{3,w}=
 \begin{cases}
 0,&w\ \text{even},\\
 \lfloor(w+1)/6\rfloor,&w\ \text{odd},
 \end{cases}
 \qquad
 \sum_{w\geq2}\dim V_{3,w}\,t^w
       =\frac{t^5}{(1-t^2)(1-t^6)}.
\]
Writing \(\Delta=XY(X-Y)\), \(Q=X^2-XY+Y^2\), the dual is parameterized
by the homogeneous polynomials
\[
 A(X,Y)=\Delta\,P(Q,\Delta^2)
\]
of degree \(w-2\).
\end{theorem}

\begin{proof}
Up to conjugation the four pairs are
\((0,1),(1,0),(1,1),(1,2)\).
As before the first two vanish and the shuffle with a zeta factor kills
the inverse-argument polynomial. The remaining polynomial \(A\)
satisfies exactly
\[
 A(Y,X)=-A(X,Y),\qquad A(X,X-Y)=A(X,Y).
\]
Antisymmetry implies divisibility by \(X-Y\); the second symmetry then
forces divisibility by \(Y\), and antisymmetry forces divisibility by
\(X\). Hence \(A=\Delta B\). The polynomial \(\Delta\) has the same
two transformation signs as \(A\), so \(B\) is invariant under swapping
\(X,Y\) and under \((X,Y)\mapsto(X,X-Y)\).

On the plane \(a+b+c=0\), with
\((a,b,c)=(X,-Y,Y-X)\), the second transformation swaps \(b,c\);
the first is global negation followed by the transposition of \(a,b\).
The generated group contains global negation, obtained by cubing
the product, and all permutations. The elementary symmetric-polynomial
theorem therefore identifies its invariant ring with
\[
 \Q[e_2,e_3^2]=\Q[Q,\Delta^2].
\]
Conversely every \(\Delta P(Q,\Delta^2)\) has the required symmetries,
and direct substitution verifies all the original color rows.
Counting monomials of degrees \(3+2a+6b=w-2\) gives the result.
Explicitly, at odd \(w\geq5\) a basis is
\[
 \Delta^{2b+1}Q^{(w-5)/2-3b},
 \qquad 0\leq b\leq\lfloor(w-5)/6\rfloor .
\]
\end{proof}

For comparison, the exact dimensions at the first few odd weights are
\[
\begin{array}{c|rrrrrrr}
 w&3&5&7&9&11&13&15\\ \hline
 \dim V_{3,w}&0&1&1&1&2&2&2\\
 \dim V_{4,w}&1&2&3&4&5&6&7
\end{array}
\]
and all even-weight dimensions are zero. The exact matrix script
checks both formulas through weight 16, including that the displayed
polynomial bases annihilate every raw row and have the stated rank.

\subsection{A separating functional and the scope of the S4 proof}

At every odd weight \(w\geq3\), choose
\(A(X,Y)=(Y-X)^{w-2}\) in the level-four theorem.
Then \(K(X,Y)=X^{w-2}\), \(V(X,Y)=-Y^{w-2}\).
This gives a rational linear functional \(\ell\) that annihilates every
defining row and every background value:
\begin{align}
 \ell(D_a(1,1))&=(-1)^{a+1}\binom{w-2}{a-1},\notag\\
 \ell(D_{w-1}(1,2))&=1,\qquad
 \ell(D_1(2,1))=-1. \label{cycloquot:eq:separating}
\end{align}
All other entries vanish, apart from conjugates with opposite signs.

At weight five its nonzero entries are particularly transparent:
\[
\begin{array}{c|rrrrrr}
 \text{symbol}&A_{1,4}&A_{2,3}&A_{3,2}&A_{4,1}
       &\Imag\Li_{1,4}(-1,\iu)&K_{4,1}\\ \hline
 \ell&1&-3&3&-1&-1&1
\end{array}
\]
The polynomial proof verifies all weights; the delivered certificate
also checks this six-entry vector against every exact weight-five row.

The now-proved identity of Theorem~\ref{gaussian:thm:S4} has an
informative image in this restricted presentation.

The exact identity \eqref{cycloquot:eq:K41} makes this equivalent to
\begin{equation}\label{cycloquot:eq:S4-A}
 161280A_{1,4}+69120(g_{4,1}+g_{3,2}+3g_{2,3})
       +617\pi^5-101520G\zeta(3)=0.
\end{equation}
The functional \(\ell\) sends its left side to \(161280\).
Consequently this relation does not follow from
\eqref{cycloquot:eq:formal-stuffle}--\eqref{cycloquot:eq:formal-shuffle} and the stated
background alone. This is an exact obstruction to a specific proof
strategy. The convergent octahedral relation used in its proof supplies precisely information outside this presentation.

\subsection{A new weight-seven candidate}

\begin{proposition}[Two weight-seven coordinate baskets]
\label{s6change:prop:coordinates}
The Gaussian doubles satisfy the exact identity
\begin{align}
 g_{2,5}={}&30g_{6,1}+15g_{5,2}+5g_{4,3}
 -\frac{15}{512}G\zeta(5)\notag\\
 &+\frac3{16}\beta(4)\zeta(3)-\frac{11\pi^7}{368640}.
 \label{s6change:eq:identity}
\end{align}
\end{proposition}
\begin{proof}
The convergent same-point shuffles give
\begin{align*}
 g_{2,5}+2g_{3,4}+3g_{4,3}+5g_{5,2}+10g_{6,1}
 &=\operatorname{Im}(\Li_2(i)\Li_5(i)),\\
 g_{3,4}+4g_{4,3}+10g_{5,2}+20g_{6,1}
 &=\operatorname{Im}(\Li_3(i)\Li_4(i)).
\end{align*}
Subtract twice the second row from the first and use
\[
 \operatorname{Im}(\Li_2(i)\Li_5(i))
 =-\frac{15}{512}G\zeta(5)-\frac{5\pi^7}{73728},
\]
\[
 \operatorname{Im}(\Li_3(i)\Li_4(i))
 =-\frac3{32}\beta(4)\zeta(3)-\frac{7\pi^7}{368640}.
\]
These product evaluations follow from the ordinary depth-one fourth-root
formulas. The coefficient of $g_{3,4}$ cancels, giving the identity.
\end{proof}
The new Cayley report's proposed formula
\begin{align*}
 S_6={}&-\frac{12334}{527}g_{6,1}-\frac{384}{31}g_{5,2}
 -\frac{2136}{527}g_{4,3}+\frac{3179\pi^7}{5713920}\\
 &-\frac{801}{2108}\beta(4)\zeta(3)-2\beta(6)\log2
\end{align*}
is therefore exactly the same conjecture as the following earlier form,
written in a different Gaussian basket. If the two residuals mean left
side minus right side, their difference is $-128/155$ times the zero
relation obtained by moving \eqref{s6change:eq:identity} to the left.
The exact two-row check has zero rational residual. This proves equivalence
of the candidates, not the conjectured vanishing of either one.

\begin{conjecture}[A reduction for \(S_6\)]\label{cycloquot:conj:S6}
With the same definitions,
\begin{align}
 S_6={}&\frac{722}{527}g_{6,1}
       +\frac{40}{527}g_{4,3}
       -\frac{128}{155}g_{2,5}
       +\frac{15191}{28569600}\pi^7\notag\\
      &-\frac3{124}G\zeta(5)
       -\frac{2373}{10540}\beta(4)\zeta(3)
       -2\beta(6)\log2. \label{cycloquot:eq:S6}
\end{align}
\end{conjecture}

The candidate was found in the ordered basket
\[
 (S_6,g_{6,1},g_{4,3},g_{2,5},\pi^7,
          G\zeta(5),\beta(4)\zeta(3),\beta(6)\log2).
\]
Its primitive integer relation vector is
\[
 \begin{split}
 (&485683200,-665395200,-36864000,401080320,\\
  &-258247,11750400,109347840,971366400).
 \end{split}
\]
A 100-digit search used coefficient bound \(10^{12}\) and tolerance
\(10^{-90}\). After the coefficients were frozen, reevaluation by
Mellin integrals at 220 working digits gave an integer-normalized
residual of approximately \(-2.315\times10^{-213}\).
An independent Euler transformation of the defining alternating
series, using 900 terms and 300 working digits, gave approximately
\(2.854\times10^{-263}\). The script records the full values and
the exact integer vector. These are numerical evidence, not a proof
or an interval certificate.

For the second method, the defining series are
\[
 g_{a,b}=\sum_{n\geq0}
       \frac{(-1)^nH_{2n}^{(b)}}{(2n+1)^a},
 \qquad H_m^{(b)}=\sum_{j=1}^{m}j^{-b},
\]
and the analogous series with \(H_n\) for \(S_6\).
Thus the second evaluation does not reuse the discovery integral.
Last transformed terms are recorded as diagnostics; their small
size alone is not asserted to bound the complete tail.

The weight-seven functional \eqref{cycloquot:eq:separating} sends the
integer-normalized candidate to \(485683200\). Its proof therefore
also requires information beyond the restricted single-product
presentation. The S4 certificate suggests seeking a corresponding weight-seven functional equation
for \eqref{cycloquot:eq:S6}, with the
nonzero functional serving as a check that the missing relation has
actually been introduced.

